\documentclass[11pt,a4paper]{amsart}
\usepackage[margin=1in]{geometry}
\usepackage{amsmath,amssymb,amsthm,mathtools}
\usepackage[hidelinks]{hyperref}
\usepackage{enumitem}

\newtheorem{theorem}{Theorem}[section]
\newtheorem{lemma}[theorem]{Lemma}
\newtheorem{proposition}[theorem]{Proposition}
\newtheorem{corollary}[theorem]{Corollary}
\theoremstyle{definition}
\newtheorem{definition}[theorem]{Definition}
\newtheorem{remark}[theorem]{Remark}

\newcommand{\wone}{\omega_1}
\newcommand{\N}{\mathbb N}
\newcommand{\otp}{\operatorname{otp}}
\newcommand{\cf}{\operatorname{cf}}

\title[An ordinal Ramsey counterexample]{An ordinal Ramsey counterexample at $\omega_1^2$}
\author{}
\date{July 2026}

\begin{document}

\begin{abstract}
We answer negatively the following infinite variant of an ordinal Ramsey question.  There is a countable graph $G$ with no $K_4$ and no $K_{\aleph_0,\aleph_0}$ such that
\[
        \omega_1^2 \nrightarrow (\omega_1\omega,G)^2.
\]
The graph $G$ may be taken to be the countable shift graph.  The coloring witnessing the negative relation is explicit: identify $\omega_1^2$ with $\omega_1\times\omega_1$ in lexicographic order and color a pair blue exactly when the two coordinates move in opposite directions.  We also record the finite case: by the finite Erd\H{o}s--Milner theorem,
\[
        \omega_1^2 \to (\omega_1\omega,H)^2
\]
for every finite graph $H$.
\end{abstract}
\maketitle

\section{Conventions and statement of the result}

All graphs in this note are simple and undirected.  Copies of graphs are not required to be induced.  For an ordinal $\alpha$ and a graph $G$, the notation
\[
        \alpha\to(\beta,G)^2
\]
means that every red/blue coloring of $[\alpha]^2$ contains either a red-homogeneous subset of order type $\beta$, or a blue copy of $G$.

Ordinal multiplication is used throughout.  Thus
\[
        \omega_1^2=\omega_1\cdot\omega_1,
        \qquad
        \omega_1\omega=\omega_1\cdot\omega.
\]
We identify $\omega_1^2$ with $\omega_1\times\omega_1$ ordered lexicographically:
\[
        (\xi,\eta)<_{\rm lex}(\zeta,\theta)
        \quad\Longleftrightarrow\quad
        \xi<\zeta\text{ or }(\xi=\zeta\text{ and }\eta<\theta).
\]
The set $\{\xi\}\times\omega_1$ is called the $\xi$-th vertical block.

The main theorem is the following.

\begin{theorem}\label{thm:main-negative}
There is a countable graph $G$ containing neither $K_4$ nor $K_{\aleph_0,\aleph_0}$ such that
\[
        \omega_1^2\nrightarrow(\omega_1\omega,G)^2.
\]
Indeed, $G$ may be chosen to be the countable shift graph.
\end{theorem}

The finite case is different.

\begin{theorem}\label{thm:finite-case}
For every finite graph $H$,
\[
        \omega_1^2\to(\omega_1\omega,H)^2.
\]
In particular, the positive relation holds for every finite graph $H$ with no $K_4$.
\end{theorem}

The proof of Theorem~\ref{thm:main-negative} is completely explicit and self-contained.  The proof of Theorem~\ref{thm:finite-case} is an immediate application of the classical finite Erd\H{o}s--Milner theorem; see Section~\ref{sec:finite-case}.

\section{The coloring of $\omega_1^2$}

Let
\[
        X=\omega_1\times\omega_1
\]
with the lexicographic order.  Define a red/blue coloring $c$ of $[X]^2$ as follows.  For distinct vertices
\[
        x=(\xi,\eta),\qquad y=(\zeta,\theta),
\]
color $\{x,y\}$ blue if, after possibly interchanging $x$ and $y$,
\[
        \xi<\zeta\quad\text{and}\quad \eta>\theta.
\]
Equivalently, a pair is blue exactly when the first coordinates and second coordinates are ordered in opposite directions.  All remaining pairs are colored red.

We first show that this coloring has no red copy of $\omega_1\omega$.

\begin{lemma}\label{lem:sections-small-order-type}
Let $A\subseteq\omega_1\times\omega_1$.  For $\xi<\omega_1$, put
\[
        A_\xi=\{\eta<\omega_1:(\xi,\eta)\in A\}.
\]
If $A_\xi$ is countable for all but at most one $\xi<\omega_1$, then $A$, with the inherited lexicographic order, has order type at most $\omega_1\cdot2$.
\end{lemma}

\begin{proof}
The inherited lexicographic order on $A$ is the ordinal sum
\[
        \sum_{\xi<\omega_1}\otp(A_\xi),
\]
where empty sections contribute $0$.  If every section is countable, then every proper initial partial sum is countable: indeed, before a fixed $\xi<\omega_1$ there are only countably many earlier indices, and a countable ordinal sum of countable ordinals is countable.  Hence the full sum is at most $\omega_1$.

Now suppose exactly one section, say $A_{\xi_0}$, is uncountable.  The part before $\xi_0$ is countable, because $\xi_0$ itself is a countable ordinal.  The section $A_{\xi_0}$ has order type at most $\omega_1$.  The part after $\xi_0$ is an ordinal sum of $\omega_1$ many countable ordinals, and therefore has order type at most $\omega_1$.  Thus the whole order type is at most
\[
        \text{countable}+\omega_1+\omega_1=\omega_1\cdot2.
\]
This proves the claim.
\end{proof}

\begin{lemma}\label{lem:no-red-omega1omega}
The coloring $c$ has no red-homogeneous subset of order type $\omega_1\omega$.
\end{lemma}

\begin{proof}
Let $A\subseteq X$ be red-homogeneous.  We claim that $A$ has at most one uncountable vertical section.

Suppose otherwise.  Choose $\xi<\zeta<\omega_1$ such that both
\[
        A_\xi=\{\eta:(\xi,\eta)\in A\},
        \qquad
        A_\zeta=\{\theta:(\zeta,
        \theta)\in A\}
\]
are uncountable.  An uncountable subset of $\omega_1$ is unbounded in $\omega_1$.  Choose any $\theta\in A_\zeta$.  Since $A_\xi$ is unbounded, choose $\eta\in A_\xi$ with $\eta>\theta$.  Then
\[
        \xi<\zeta,
        \qquad
        \eta>\theta,
\]
so the pair $\{(\xi,\eta),(\zeta,\theta)\}$ is blue by definition of $c$, contradicting the red-homogeneity of $A$.

Therefore $A$ has at most one uncountable vertical section.  By Lemma~\ref{lem:sections-small-order-type}, $A$ has order type at most $\omega_1\cdot2$.  Since
\[
        \omega_1\cdot2 < \omega_1\cdot\omega,
\]
$A$ cannot have order type $\omega_1\omega$.
\end{proof}

\section{The forbidden graph: the countable shift graph}

\begin{definition}
The countable shift graph $S_\omega$ has vertex set
\[
        V(S_\omega)=\{(i,j): i<j<\omega\}.
\]
Two distinct vertices $(i,j)$ and $(k,\ell)$ are adjacent if and only if, after possibly interchanging them,
\[
        j=k.
\]
Equivalently, the edges are exactly
\[
        (i,j)\sim(j,k)
        \qquad(i<j<k<\omega).
\]
\end{definition}

We record the two structural properties needed for the Ramsey counterexample.

\begin{lemma}\label{lem:shift-triangle-free}
The shift graph $S_\omega$ is triangle-free.  In particular, it contains no $K_4$.
\end{lemma}

\begin{proof}
Fix a vertex $v=(i,j)$.  Its neighbors are precisely
\[
        L_v=\{(h,i):h<i\}
        \qquad\text{and}\qquad
        R_v=\{(j,k):j<k<\omega\}.
\]
Two vertices of $L_v$ are not adjacent: if $h<h'<i$, then adjacency of $(h,i)$ and $(h',i)$ would require $i=h'$, impossible because $h'<i$.  Two vertices of $R_v$ are not adjacent: if $j<k<k'$, then adjacency of $(j,k)$ and $(j,k')$ would require $k=j$, impossible.  Finally, a vertex $(h,i)\in L_v$ and a vertex $(j,k)\in R_v$ are not adjacent, since adjacency would require either $i=j$ or $k=h$, while
\[
        h<i<j<k.
\]
Thus the neighborhood of every vertex is independent.  Hence no vertex can lie in a triangle, and $S_\omega$ is triangle-free.
\end{proof}

\begin{lemma}\label{lem:shift-no-biclique}
The shift graph $S_\omega$ contains no copy of $K_{\aleph_0,\aleph_0}$.
\end{lemma}

\begin{proof}
We first observe that two vertices of $S_\omega$ with distinct second coordinates have only finitely many common neighbors.

Let
\[
        u=(a,b),\qquad v=(c,d),
\]
where $a<b<\omega$, $c<d<\omega$, and $b\ne d$.  A vertex $w=(x,y)$ is adjacent to $u$ precisely when
\[
        y=a\quad\text{or}\quad x=b,
\]
and it is adjacent to $v$ precisely when
\[
        y=c\quad\text{or}\quad x=d.
\]
Thus a common neighbor must satisfy one of the following four systems:
\[
\begin{array}{ll}
        y=a\text{ and }y=c, & y=a\text{ and }x=d,\\
        x=b\text{ and }y=c, & x=b\text{ and }x=d.
\end{array}
\]
The fourth system is impossible because $b\ne d$.  The second and third systems determine at most one vertex each.  The first system is impossible unless $a=c$, and if $a=c$ then it permits only vertices $(x,a)$ with $x<a$, of which there are finitely many.  Hence $u$ and $v$ have finitely many common neighbors.

Now suppose toward a contradiction that $S_\omega$ contains a $K_{\aleph_0,\aleph_0}$ with bipartition $(A,B)$.  The side $A$ is infinite.  For any fixed $n<\omega$, there are only finitely many vertices of $S_\omega$ with second coordinate $n$, namely $(0,n),\dots,(n-1,n)$.  Therefore $A$ contains two vertices with distinct second coordinates.  These two vertices have only finitely many common neighbors by the preceding paragraph.  But every vertex of $B$ is a common neighbor of them, and $B$ is infinite.  This contradiction proves the lemma.
\end{proof}

It remains to show that the coloring from Lemma~\ref{lem:no-red-omega1omega} has no blue copy of $S_\omega$.

\begin{lemma}[Infinite Ramsey theorem, finite-uniform form]\label{lem:ordinary-ramsey}
For every $m<\omega$ and every finite coloring of $[\omega]^m$, there is an infinite subset of $\omega$ whose $m$-element subsets all have the same color.
\end{lemma}

\begin{proof}
We prove the result by induction on $m$.  The cases $m=0$ and $m=1$ are immediate from the infinite pigeonhole principle.  Assume the theorem holds for $m$ and consider a finite coloring $d$ of $[\omega]^{m+1}$.

Set $M_{-1}=\omega$.  Construct recursively an increasing sequence
\[
        n_0<n_1<n_2<\cdots
\]
and infinite sets
\[
        M_0\supseteq M_1\supseteq M_2\supseteq\cdots
\]
with $n_s<\min M_s$ as follows.  Having chosen $M_{s-1}$, let $n_s=\min M_{s-1}$ and consider the induced coloring of $[M_{s-1}\setminus(n_s+1)]^m$ given by
\[
        F\longmapsto d(\{n_s\}\cup F).
\]
By the induction hypothesis, there is an infinite $M_s\subseteq M_{s-1}\setminus(n_s+1)$ on which this induced coloring is constant.  Denote the resulting color by $e(s)$.

Since there are only finitely many colors, choose an infinite set $S\subseteq\omega$ such that $e(s)$ is constant for $s\in S$.  Then $\{n_s:s\in S\}$ is homogeneous for $d$: if
\[
        s_0<s_1<\cdots<s_m
\]
are in $S$, then $\{n_{s_1},\dots,n_{s_m}\}\subseteq M_{s_0}$, and by the construction at stage $s_0$ the color of
\[
        \{n_{s_0},n_{s_1},\dots,n_{s_m}\}
\]
is $e(s_0)$, independent of the chosen $(m+1)$-set.  This completes the induction.
\end{proof}

\begin{lemma}\label{lem:no-blue-shift}
The blue graph of the coloring $c$ contains no copy of $S_\omega$.
\end{lemma}

\begin{proof}
Suppose toward a contradiction that there is a blue embedding
\[
        \varphi:S_\omega\hookrightarrow X.
\]
Write
\[
        \varphi(i,j)=(\xi_{ij},\eta_{ij})
        \qquad(i<j<\omega).
\]
For every triple $i<j<k<\omega$, the vertices $(i,j)$ and $(j,k)$ are adjacent in $S_\omega$.  Hence their images are blue-adjacent in $X$.  Therefore exactly one of the following two alternatives holds:
\[
\begin{array}{ll}
        \text{type }0: & \xi_{ij}<\xi_{jk}\text{ and }\eta_{ij}>\eta_{jk},\\[2mm]
        \text{type }1: & \xi_{jk}<\xi_{ij}\text{ and }\eta_{jk}>\eta_{ij}.
\end{array}
\]
Color the triple $\{i,j,k\}$ by its type.  By Lemma~\ref{lem:ordinary-ramsey}, there is an infinite set
\[
        M=\{m_0<m_1<m_2<\cdots\}\subseteq\omega
\]
whose triples all have the same type.

If all triples from $M$ have type $0$, then for each $t<\omega$ the triple
\[
        m_t<m_{t+1}<m_{t+2}
\]
gives
\[
        \eta_{m_t m_{t+1}}>\eta_{m_{t+1}m_{t+2}}.
\]
Hence
\[
        \eta_{m_0m_1}>\eta_{m_1m_2}>\eta_{m_2m_3}>\cdots
\]
is an infinite strictly decreasing sequence of ordinals, impossible.

If all triples from $M$ have type $1$, then similarly
\[
        \xi_{m_0m_1}>\xi_{m_1m_2}>\xi_{m_2m_3}>\cdots
\]
is an infinite strictly decreasing sequence of ordinals, again impossible.

Both cases are impossible.  Therefore no blue copy of $S_\omega$ exists.
\end{proof}

\begin{proof}[Proof of Theorem~\ref{thm:main-negative}]
Let $G=S_\omega$.  By Lemma~\ref{lem:shift-triangle-free}, $G$ contains no $K_4$.  By Lemma~\ref{lem:shift-no-biclique}, $G$ contains no $K_{\aleph_0,\aleph_0}$.  The coloring $c$ of $[\omega_1^2]^2$ has no red copy of $\omega_1\omega$ by Lemma~\ref{lem:no-red-omega1omega}, and no blue copy of $G$ by Lemma~\ref{lem:no-blue-shift}.  Hence
\[
        \omega_1^2\nrightarrow(\omega_1\omega,G)^2.
\]
\end{proof}

\section{The finite case}\label{sec:finite-case}

The infinite counterexample above cannot be finite.  The relevant positive theorem is the following finite form of the Erd\H{o}s--Milner ordinal Ramsey theorem.

\begin{theorem}[Finite Erd\H{o}s--Milner theorem]\label{thm:finite-erdos-milner}
For every finite integer $r\ge1$,
\[
        \omega_1^2\to(\omega_1\omega,r)^2.
\]
Equivalently, every red/blue coloring of $[\omega_1^2]^2$ contains either a red-homogeneous subset of order type $\omega_1\omega$, or a blue complete graph on $r$ vertices.
\end{theorem}

\begin{remark}
Theorem~\ref{thm:finite-erdos-milner} is a classical theorem of ordinal partition calculus.  It is usually quoted as a finite Erd\H{o}s--Milner theorem.  The present note uses it only in this standard finite form.  The negative construction in Theorem~\ref{thm:main-negative} is independent of this theorem.
\end{remark}

\begin{proof}[Proof of Theorem~\ref{thm:finite-case}]
Let $H$ be a finite graph and put $r=|V(H)|$.  Given a red/blue coloring of $[\omega_1^2]^2$, Theorem~\ref{thm:finite-erdos-milner} gives either a red-homogeneous subset of order type $\omega_1\omega$, or a blue $K_r$.  In the second case, because graph copies are not required to be induced, a blue $K_r$ contains a blue copy of every graph on $r$ vertices, and therefore contains a blue copy of $H$.  Thus
\[
        \omega_1^2\to(\omega_1\omega,H)^2
\]
for every finite graph $H$.
\end{proof}

\section{Conclusion}

The proposed infinite assertion is false.  The countable shift graph $S_\omega$ satisfies
\[
        K_4\nsubseteq S_\omega,
        \qquad
        K_{\aleph_0,\aleph_0}\nsubseteq S_\omega,
\]
but nevertheless
\[
        \omega_1^2\nrightarrow(\omega_1\omega,S_\omega)^2.
\]
The witnessing coloring is explicit and uses the inversion relation on $\omega_1\times\omega_1$.

For finite graphs the answer is positive: every finite graph $H$ satisfies
\[
        \omega_1^2\to(\omega_1\omega,H)^2.
\]
Thus the obstruction is genuinely infinite, even though the counterexample graph is only countable and is sparse enough to avoid both $K_4$ and $K_{\aleph_0,\aleph_0}$.

\begin{thebibliography}{9}

\bibitem{ErdosMilner}
P. Erd\H{o}s and E. C. Milner,
\emph{A theorem in the partition calculus},
Canadian Mathematical Bulletin \textbf{15} (1972), 501--505.

\end{thebibliography}

\end{document}
